\documentclass[a4paper]{article}
\usepackage[pdftex]{graphicx}
\usepackage[utf8]{inputenc}
\usepackage{enumerate}
\usepackage{siunitx}
\sisetup{locale=DE} 
\usepackage{href-ul}
\hypersetup{
	colorlinks=true,
	linkcolor=blue,
	urlcolor=blue}
\usepackage{icomma}
\usepackage{amssymb}
\usepackage{geometry}
\geometry{a4paper, top=15mm, left=15mm, right=15mm, bottom=15mm,
	headsep=10mm, footskip=12mm}


% Start the document
\begin{document}
%Titel
{\bf Schulaufgabe aus der Physik}\\
\begin{enumerate}[1.]

{\bf \item Auf einen Doppelspalt fällt Laserlicht der Wellenlänge $\lambda=\SI{546}{\nano\meter}$.} Auf einem 2,80 m entfernten Schirm beobachtet man helle Streifen. 
\begin{enumerate}[a)]
\item Auf dem Schirm haben die beiden Maxima 1. Ordnung einen Abstand von 17,6 cm.\\
Berechnen Sie den Spaltabstand.
\item Wie ändern sich Anzahl und Lage der Maxima, wenn bei ansonsten unveränderter Anordnung Laserlicht einer kleineren Wellenlänge verwendet wird?
\end{enumerate}

{\bf \item Wie lang müsste der Schirmabstand} mindestens sein, um die Natrium-Doppellinie ($\SI{588,99}{\nano\meter}$ und $\SI{589,59}{\nano\meter}$)mit einem Gitter der Spaltabstände $b=\SI{100}{\micro\meter}$ aufzulösen, sodass der Abstand beider Maxima 1. Ordnung $\SI{3}{\milli\meter}$ ist?


\end{enumerate} 

\fbox{
	\begin{minipage}{0.5\textwidth}
		Zur Lösung bitte \href{https://www.okuyakl.de/phys/p11wellL095/ll095.pdf}{hier klicken} oder den QR-Code scannen.\\
	Weitere Arbeitsblätter gibt es unter 
	
	\href{https://www.okuyakl.de}{www.okuyakl.de}
	\end{minipage}
	\hfill
	\begin{minipage}{0.4\textwidth}
		\includegraphics[width=1.5 cm]{../../viecher/zwe03}
		\includegraphics[width=3 cm]{qr095}
		\includegraphics[width=2 cm]{../../viecher/afanticon1}
		
	\end{minipage}}

\end{document}%Lösung---------------------------------------
